\subsection{Existence and uniqueness of the holomorphic functional calculus}

Suppose that A is a commutative unital complex Banach algebra.

DEFINITION 2. --- Let \(n \geqslant 1\) be an integer and \(\boldsymbol{a} \in \mathrm{A}^{n}\) . Let U be an open neighborhood of \(\mathrm{Sp}^{n}(\boldsymbol{a})\). We say that a family \(\boldsymbol{a}'\) is an envelope of \((\boldsymbol{a}, \mathrm{U})\) if \(\boldsymbol{a}' \in \mathbf{C}^{n+p}\) extends \(\boldsymbol{a}\) and if \(\mathrm{U} \times \mathbf{C}^{p}\) contains the polynomially convex envelope of \(\mathrm{Sp}^{n+p}(\boldsymbol{a}')\) .

Lemma 11. --- Let \(n \geqslant 1\) be an integer. Let \(\boldsymbol{a} \in \mathrm{A}^{n}\) . For every open neighborhood U of \(\mathrm{Sp}^{n}(\boldsymbol{a})\) , there exists an envelope of \((\boldsymbol{a}, \mathrm{U})\) .

Let \((a_{\lambda})_{\lambda\in\Lambda}\) be a family of elements of A extending the family a and topologically generating the unital Banach algebra A. Let \(\pi\) be the canonical projection from \(C^{\Lambda}\) onto \(C^{n}\) and let \(U'=\pi^{-1}(U)\) . Then \(U'\) is a neighborhood of \(\mathrm{Sp}^{\Lambda}((a_{\lambda}))\) , and \(\mathrm{Sp}^{\Lambda}((a_{\lambda}))\) is polynomially convex (I, p.~44, lemma 4). By lemma 6 of I, p.~47, there exists a finite subset \(\Lambda_{0}\) of \(\Lambda\) containing \(\{1,2,\ldots,n\}\) such that \(\mathrm{pr}_{\Lambda_{0}}(\mathrm{U}')\) contains the polynomially convex hull S of \(\mathrm{pr}_{\Lambda_{0}}(\mathrm{Sp}^{\Lambda}((a_{\lambda})_{\lambda\in\Lambda}))=\mathrm{Sp}^{\Lambda_{0}}((a_{\lambda})_{\lambda\in\Lambda_{0}})\) . Let \(p\geqslant0\) be the integer such that \(\Lambda_{0}\) has cardinality \(n+p\) and let j be a bijection from \(\{1,\ldots,n+p\}\) into \(\Lambda_{0}\) which coincides with the identity map on \(\{1,\ldots,n\}\). Since the projection of S is contained in U, the family \((a_{j(k)})_{1\leqslant k\leqslant n+p}\) is an envelope of \((\boldsymbol{a},\mathrm{U})\) .

PROPOSITION 5. --- The data of the mappings \(\Theta_{\boldsymbol{a}}\) , for \(n \geqslant 1\) and \(\boldsymbol{a} \in \mathrm{A}^{n}\) , is a holomorphic functional calculus on A, that is, conditions (CF1), (CF2), and (CF3) of I, p.~51 are satisfied.

Let \(n \geqslant 1\) be an integer and \(\boldsymbol{a} = (a_{1}, \ldots, a_{n}) \in \mathrm{A}^{n}\) . The map \(\Theta_{\boldsymbol{a}}\) satisfies \(\Theta_{\boldsymbol{a}}(z_{i}) = a_{i}\) for every i such that \(1 \leqslant i \leqslant n\) by lemma 10 of I, p.~64, which proves property (CF2). Lemma 9 of I, p.~63 implies property (CF3) for the maps \(\Theta_{\boldsymbol{a}}\) .

The map \(\Theta_{\pmb{a}}\) is A-linear and continuous (I, p.~61, no. 5). It satisfies \(\Theta_{\pmb{a}}(1) = 1\) (lemma 9 of I, p.~63). To verify condition (CF1), it remains to establish that \(\Theta_{\pmb{a}}\) is an algebra morphism. To do this, we will prove that \(\Theta_{\pmb{a}}^{\mathrm{U}}\) is an algebra morphism for every open neighborhood U of \(\mathrm{Sp}^{n}(\pmb{a})\) .

Suppose first that U contains the polynomially convex hull K of \(\mathrm{Sp}^{n}(\boldsymbol{a})\) . Let \(f_{1}\) and \(f_{2}\) be elements of \(\mathcal{O}(\mathrm{U};\mathrm{A})\) . There exists a sequence \((f_{1,k})\) (resp. \((f_{2,k})\)) of polynomial functions which converges to \(f_{1}\) (resp. to \(f_{2}\) ) in \(\mathcal{O}(\mathrm{K};\mathrm{A})\) (Theorem 2 of I, p.~68), hence in \(\mathcal{O}(\mathrm{Sp}^{n}(\boldsymbol{a});\mathrm{A})\) . For every integer \(k\) , one has

\[
\Theta_ {\boldsymbol {a}} ^ {\mathrm{U}} (f _ {1, k}) \Theta_ {\boldsymbol {a}} ^ {\mathrm{U}} (f _ {2, k}) = \Theta_ {\boldsymbol {a}} ^ {\mathrm{U}} (f _ {1, k} f _ {2, k})
\]

by Lemma 10 of I, p.~64, whence \(\Theta_{a}^{\mathrm{U}}(f_{1})\Theta_{a}^{\mathrm{U}}(f_{2}) = \Theta_{a}^{\mathrm{U}}(f_{1}f_{2})\) by passing to the limit.

Consider the general case. Let \(a' \in C^{n+p}\) be an envelope of \((\boldsymbol{a}, U)\) (lemma 11) and \(\pi: U \times C^{p} \to U\) the canonical projection. Since \(U \times C^{p}\) contains the polynomially convex envelope of \(\mathrm{Sp}^{n+p}(\boldsymbol{a}')\) , one has

\[
\Theta_ {\boldsymbol {a} ^ {\prime}} ^ {\mathrm{U} \times \mathbf {C} ^ {p}} (f _ {1} \circ \pi) \Theta_ {\boldsymbol {a} ^ {\prime}} ^ {\mathrm{U} \times \mathbf {C} ^ {p}} (f _ {2} \circ \pi) = \Theta_ {\boldsymbol {a} ^ {\prime}} ^ {\mathrm{U} \times \mathbf {C} ^ {p}} (f _ {1} f _ {2} \circ \pi)
\]

for \(f_1\) and \(f_2\) in \(\mathcal{O}(\mathrm{U};\mathrm{A})\) by the first case. As, for every function \(f \in \mathcal{O}(\mathrm{U};\mathrm{A})\) , one has \(\Theta_{\boldsymbol{a}'}^{\mathrm{U}\times \mathbf{C}^p}(f\circ \pi) = \Theta_{\boldsymbol{a}}^{\mathrm{U}}(f)\) (condition (CF3) previously proved), the conclusion follows, and hence condition (CF1).

We can now prove Theorem 1 of I, p.~51. Prop. 5 shows that the family of maps \((\Theta_{a})_{a}\) is a holomorphic functional calculus on A. It therefore remains only to establish the uniqueness of the holomorphic functional calculus on A.

Let \((\Psi_{a})_{a}\) be a family of maps defined for every integer \(n \geqslant 1\) and every \(a \in A^{n}\) and satisfying conditions (CF1), (CF2), (CF3) of the holomorphic functional calculus on A (I, p.~51). It suffices to prove that for every integer \(n \geqslant 1\), every \(a \in A^{n}\), and every open neighborhood U of a, we have \(\Theta_{a}^{U} = \Psi_{a}^{U}\) .

Let \(n \geqslant 1\) and \(\boldsymbol{a} = (a_{1}, \ldots, a_{n}) \in \mathrm{A}^{n}\) . Let U be an open neighborhood of \(\mathrm{Sp}^{n}(\boldsymbol{a})\) . Suppose first that U contains the polynomially convex hull K of \(\mathrm{Sp}^{n}(\boldsymbol{a})\) . The morphisms \(\Theta_{a}^{U}\) and \(\Psi_{a}^{U}\) coincide on polynomial functions by properties (CF1) and (CF2). By the corollary of Theorem 2 of I, p.~68 and the continuity property (CF1), these morphisms are therefore equal.

Let us prove the general case. Let \(a' \in C^{n+p}\) be an envelope of \((\boldsymbol{a}, U)\) and \(\pi: U \times C^{p} \to U\) the canonical projection. We have

\[
\Theta_ {\boldsymbol {a}} ^ {\mathrm{U}} = \Theta_ {\boldsymbol {a} ^ {\prime}} ^ {\mathrm{U} \times \mathbf {C} ^ {p}} \circ \pi^ {*} = \Psi_ {\boldsymbol {a} ^ {\prime}} ^ {\mathrm{U} \times \mathbf {C} ^ {p}} \circ \pi^ {*} = \Psi_ {\boldsymbol {a}} ^ {\mathrm{U}}
\]

by property (CF3) and the preceding case. This concludes the proof of Theorem 1 of I, p.~51.

Note that Th. 2 of I, p.~68 also entails the following uniqueness result:

PROPOSITION 6. --- Let \(\boldsymbol{a} \in \mathrm{A}^n\) . We suppose \(\mathrm{Sp}^n(\boldsymbol{a})\) polynomially convex. Let \(z_1, \ldots, z_n\) be the germs in a neighbourhood of \(\mathrm{Sp}^n(\boldsymbol{a})\) of the coordinate functions on \(\mathbf{C}^n\). Then the map \(\Theta_{\boldsymbol{a}}\) is the unique continuous morphism of unital algebras \(\varphi\) of \(\mathcal{O}(\mathrm{Sp}^n(\boldsymbol{a}); \mathrm{A})\) to A such that \(\varphi(z_1) = a_1, \ldots, \varphi(z_n) = a_n\) .

Lemma 10 of I, p.~64 and the corollary of Proposition 2 of I, p.~65 justify the following notations for the holomorphic functional calculus. Let \(n \geqslant 1\) and \(\boldsymbol{a} \in \mathrm{A}^n\) . For every germ \(f \in \mathcal{O}(\mathrm{K};\mathrm{A})\) (resp. for every holomorphic function \(f \in \mathcal{O}(\mathrm{U};\mathrm{A})\) on an open neighbourhood U of \(\mathrm{Sp}^n (\boldsymbol{a})\) ), we set

\[
f (\boldsymbol {a}) = \Theta_ {\boldsymbol {a}} (f).\tag{6}
\]

This notation is consistent with the notation introduced in A, IV, p.~4, n°3, if f is a polynomial, by properties (CF1) and (CF2). Properties (CF2) and (CF3) of I, p.~51 can then be written

\[
z _ {i} (\boldsymbol {a}) = a _ {i}, \quad 1 \leqslant i \leqslant n, \quad (f \circ \pi_ {m, n}) (\boldsymbol {a}) = f (\pi_ {m, n} (\boldsymbol {a})).
\]

